\chapter{The Strain and the Taper}
% OUTLINE Ch2 -- the device's dissipation SHAPE, held as DESCRIBED-IMAGE per
% Kodo's built/owed line (14:07/14:26Z): the strain frame + the taper/apodization
% are DISCIPLINES the device lives by, honest as described-image and fenced-shape,
% NOT exhibited as built fluid theorems (no standalone theorem carries them; the
% BUILT exhibit -- the period-3 residue -- is Ch3's). The taper genuinely damps
% the Gibbs ring to the floor (honestly dissipative, more honest than a box-floor
% apparent-viscosity image), but it is the model's strain structure, image-not-
% identity, never the NS PDE. CITATIONS per Kodo RULING 1 (14:26Z): the general
% dissipation function as the taper's damping shape = Rayleigh 1877; a viscous NS
% term = Stokes 1845; NOT Onsager. Gibbs phenomenon = Wilbraham 1848 / Gibbs 1899
% (inventor-first). Unabashed math + four riders; bold-math/fenced-referent; no
% Lean identifiers; step-not-move; lab-gap where any number appears; model-not-
% world. Gate: Kodo (earn/citation, described-image line) + Beastmaster (arrival-
% read: NOT-a-built-theorem held hardest).

The coarsening reveals its structures one at a time, and the first of them is a
strain. Before the model's fluid behavior can be read as flow, it must be read as
tension --- the resistance the argument puts up when it is pressed, and the way
that resistance discharges. This chapter reads that first structure, and it reads
it honestly about its own status, which is the honesty the whole volume turns on:
the strain and its relaxation are structures the instrument \emph{lives by} and
exhibits, disciplines carried in its own construction. They are not theorems. The
one built, machine-checked result of this face waits for the next chapter; here
the reading is of shape, and the chapter says so plainly wherever it reads.

Press the argument at a single point --- demand of it one exact value where it
would rather give a bracket --- and it does not yield that value cleanly. It
\emph{rings}. Around the forced point the count overshoots and undershoots,
oscillating past the demand and back, and the overshoot does not vanish as the
resolution rises; it narrows and stays. This is the shape long known in the
analysis of a series forced to represent a jump: the overshoot Henry Wilbraham
first noted in 1848 and J.\ Willard Gibbs analyzed in 1899, the \emph{Gibbs
phenomenon}, in which a sum of smooth pieces made to cross a discontinuity rings
at the crossing by a fixed fraction that refinement cannot remove. The
instrument's forced point rings in exactly that shape. The ringing \emph{is} the
strain: the argument's resistance to being pinned, made visible as an
oscillation the count carries around the site of the demand.

\begin{intuition}
The picture is a beam pressed at one spot: it does not simply bend to the load,
it quivers, and the quiver rings out along the beam and settles. The forced
value is the load; the Gibbs overshoot is the quiver; the strain is the beam's
resistance felt as that ringing. The picture is offered for the shape and fenced
for the mechanism: the instrument is no beam and computes no elasticity; what it
carries is a count that overshoots a forced point in the fixed fraction Gibbs
described, and ``strain'' is the name this book gives that carried overshoot, not
a claim of a material under load.
\end{intuition}

The honest response to a ringing is not a harder demand. Force the point more
sharply --- cut it off cleaner --- and the ring only grows; the hard edge is what
rang in the first place. The instrument's discipline, carried since its earlier
volumes, is the opposite: it does not cut, it \emph{tapers}. Rather than a hard
boundary at the forced point, it lets the demand fall off smoothly to the floor,
and the smooth fall-off damps the ring. This is the practice an instrument-maker
calls \emph{apodization} --- the removal of the foot, the smoothing of a hard
edge into a gentle one --- and its effect here is genuinely dissipative: the
taper does not hide the overshoot behind an averaged number, it \emph{damps} it,
carrying the ring's energy down to the floor the way a dissipation function
carries the energy of a disturbed system away. It is that general dissipation
function --- the one Lord Rayleigh set down in 1877 in \emph{The Theory of
Sound} to account for the energy a vibrating system loses --- that the taper
wears as its shape: a quadratic measure of how fast the disturbance is bled
away. The instrument's taper is honestly dissipative in that sense, and this is a
sharper and more honest reading than the older image of an ``apparent
viscosity'' read off the floor of a box; the taper does not merely look
viscous, it actively damps the strain it is given.

And now the fence that governs the whole chapter, stated at its plainest. All of
this --- the ringing strain, the damping taper --- is read as \emph{shape}. The
device carries the overshoot and carries the taper that damps it; these are
disciplines of its construction, exhibited here as what they are. They are not
presented as a built theorem of fluid behavior, because no single machine-checked
result stands under ``the strain relaxes'' or ``the taper damps'' the way one
will stand under the next chapter's reading. Where this book has a checked
theorem, it will exhibit it; here it has a lived discipline and a borrowed shape,
and it files them as exactly that. The dissipation is real \emph{in the device's
own terms} --- the taper genuinely bleeds the ring down --- and it is an image,
not an identity, of the dissipation a viscous fluid performs; the viscous term of
the Navier--Stokes energy equation, which George Stokes wrote in 1845, is its
distant namesake and not its content.

Two fences close the chapter, as in every chapter of this volume. The strain and
the taper are the \emph{model's} --- the electron model's own resistance and its
own damping, an image of dissipation and not a fluid of the world; nothing here
touches real viscosity or the equations that govern it. And where the taper reads
a floor, the number it reads is the model's own, at the model's resolution, never
the laboratory's, here as everywhere. The strain has been read, and the taper
that damps it; the next chapter reads what survives the damping --- the one
structure of this face that stands on a theorem, the behavior that emerges and
does not decay.
